\scanpage{33}\mspage{32}
\deleted{\unsure{forment un}} correspondent \unsure{:|à} un atelier $\mathrm{At}\ 1$ -- $\mathrm{At}\ 7$. Donc c'est le contenu du

\uline{Lemme} 6 \quad Soient $X,Y,\unsure{2|Z}$ avec
\[
 Z\overset{\circ}{\ll}X,\ Z\overset{\circ}{\ll}Y,\quad X\mathrel{\tikz[baseline=-.39ex,x=1em,y=1em]{\draw[line width=.05em] (0,.55)--(.70,.27)--(0,0) (.70,.55)--(0,.27)--(.70,0);}}Y
\]
alors $X=Y$.

Ça résultera du cor. plus général

\[
\left|\begin{aligned}
&\text{\uline{Cor} \quad On a, si }X\mathrel{\tikz[baseline=-.39ex,x=1em,y=1em]{\draw[line width=.05em] (0,.55)--(.70,.27)--(0,0) (.70,.55)--(0,.27)--(.70,0);}}Y\\
& Z\overset{\circ}{\ll}X,\quad Z\ll Y,\Longrightarrow X\leq\deleted{\unsure{<|\ll}}Y
\end{aligned}\right.
\]
\uline{Dém} \quad Soit $x\in\mathrm{omb}(Z)^\circ$ (ML 2). On a donc \deleted{$x\ll$ \unsure{om|ou|on}}
\[
 x\overset{\circ}{\ll}Z\overset{\circ}{\ll}X
\]
donc (lemme \unsure{2|?}, cor. 4 p.29) $x\overset{\circ}{\ll}X$, donc $x\in\mathrm{omb}(X)^\circ\cap\mathrm{omb}(Y)$, lequel est $\neq\varnothing$, donc $X\leq Y$ (lemme 4).

Il faut voir que la relation de \unsure{disjoint|disjont} $\parallel$ sur $L$ est bien la relation de disjonction au sens de l'atelier qu'on vient de construire, i.e.

\uline{Lemme} 7 \quad \deleted{\unsure{Soit|So}} a) Soient $x,y\in L$. Alors $x\mathrel{\tikz[baseline=-.39ex,x=1em,y=1em]{\draw[line width=.05em] (0,.55)--(.70,.27)--(0,0) (.70,.55)--(0,.27)--(.70,0);}}y$ \unsure{sss|ssi} $x=y$ ou $x\parallel y$.

b) On a $x\parallel y$ sss $x\mathrel{\tikz[baseline=-.39ex,x=1em,y=1em]{\draw[line width=.05em] (0,.55)--(.70,.27)--(0,0) (.70,.55)--(0,.27)--(.70,0);}}y$ et $\mathrm{Inf}^{\deleted{\unsure{L|F}}\unsure{\mathcal F|\mathscr F|F},\leq}(x,y)=\varnothing_{\unsure{F|\mathcal F}}$. (i.e. $x$ et $y$ disjoints au sens de l'atelier)

\unsure{Pour|Dans} a, il suffit de voir le cas $x\neq y$. À cause du \uline{caractère minimal} de $x,y$, on aura \inserted{$\mathrm{omb}\,x=\{x\},\ \mathrm{omb}(y)=\{y\},$} $\unsure{L|\llcorner}=\varnothing$, et le critère \deleted{est bien} par défaut $x\mathrel{\tikz[baseline=-.39ex,x=1em,y=1em]{\draw[line width=.05em] (0,.55)--(.70,.27)--(0,0) (.70,.55)--(0,.27)--(.70,0);}}y$ est bien $x\parallel y$. L'énoncé b) \unsure{suit;|suit,}
